% !PRINTED 459
\weberpage{459}

Nun ist nach (7) $\beta$ durch 3 nicht teilbar, und daher reduziert
sich die Kongruenz (5) durch Multiplikation mit $\beta$ auf:

\[
\alpha+\delta \equiv 0\ (\mathrm{mod}\ 3)
\]

oder nach (7) und (3) auf

\weq{10}{B \equiv 0\ (\mathrm{mod}\ 3).}

Ersetzen wir in (6) $\omega$ durch $\omega\pm 1$, so geht $B$ über in
$B\pm 2A$, und von den drei Zahlen $B, B+2A, B-2A$ ist eine und nur
eine durch 3 teilbar. Wir nehmen also an, es sei $B$ selbst durch 3
teilbar, dann folgt, daß von den drei Werten

\weq{11}{\gamma_2(\omega), \quad \gamma_2(\omega+1) = e^{-\frac{2\pi i}{3}}\gamma_2(\omega), \quad \gamma_2(\omega-1) = e^{\frac{2\pi i}{3}}\gamma_2(\omega)}

\gesperrt{nur der erste} der Gleichung

\weq{12}{\Phi_n(u, u) = 0}

genügt. Die Gleichung (12) und

\weq{13}{u^3-j(\omega) = 0}

haben daher \gesperrt{nur eine} gemeinsame Wurzel, und diese ist rational durch
$j(\omega)$ ausdrückbar. Damit ist bewiesen:

1. \gesperrt{$\gamma_2(\omega)$ ist eine Klasseninvariante, wenn
$\omega$ die Wurzel einer quadratischen Form ist, deren Diskriminante
und erster Koeffizient durch 3 nicht teilbar sind, während der
mittlere Koeffizient durch 3 teilbar ist.}

Wenden wir dies Ergebnis auf den Fall $n = 2$ an, so erhalten wir
zunächst aus der Gleichung (5), \S\,71:

\weq{14}{\Phi_2(u, u) = u^4-2u^3-495u^2+2^4\cdot 3^3\cdot 5^3 = 0.}

Die Gleichung

\[
y^2-Dx^2 = 8
\]

ist, da $D \equiv 0$ oder $\equiv 1\ (\mathrm{mod}\ 4)$ sein muß, nur
für drei negative Werte von $D$ lösbar, nämlich

\weq{15}{
\begin{aligned}
D &= -8, & x &= 1, & y &= 0, \\
D &= -7, & x &= 1, & y &= \pm 1, \\
D &= -4, & x &= 1, & y &= \pm 2.
\end{aligned}
}

Diesen drei Werten von $D$ entsprechend kann man für die Gleichung
(6) die folgenden wählen:

\[
\begin{aligned}
D &= -8, & \omega^2+2 &= 0, \\
D &= -7, & \omega^2+3\omega+4 &= 0, \\
D &= -4, & \omega^2+1 &= 0
\end{aligned}
\]

oder
%CONT>>
